% Section: Causal diagnostics and validation (manuscript agent owns prose).
\subsection{Physical-time windows}
For triples, the observation horizon is $t_h=f\,t_{\max}$. A system is included
at a horizon only if its event time exceeds $t_h$. Every feature is computed
exclusively from samples satisfying $t\le t_h$.

\subsection{Feature families}
The feature vector contains hierarchy statistics and derivatives, pair-energy
exchange, angular-momentum channels, torques, seam-tension summaries,
signal-complexity measures, and closest-approach diagnostics.
\todo{Add exact equations and a feature table.}

\subsection{Leakage controls}
The stored 79-column schema historically contained two future-length-dependent
columns: index 73 (window length over final recorded length) and index 72
(first hierarchy breach normalized by the full record). Both are identically
zero in every FINAL model input at training and serving time. Repository
invariance tests require that modifying or extending samples strictly after the
observation horizon leaves the causal vector bit-identical.

\subsection{Statistical protocol}
We report ROC--AUC, PR--AUC, balanced accuracy, Matthews correlation, Brier
score, and calibration error. Fixed false-alarm thresholds are selected on a
calibration cohort and evaluated on an untouched test cohort; per-horizon FAR
is reported separately from cumulative system-level false-alert probability.
Label-shuffle controls, blocked initial-condition splits, and no-refit
prospective scoring on the matched tail test leakage and extrapolation.
